% !PRINTED 694
\weberpage{694}
\begin{center}
§. 191.

\gesperrt{Sätze aus der Reihenlehre.}
\end{center}

Bei den weiteren Anwendungen der bisherigen Resultate sind einige Sätze aus der Lehre von den unendlichen Reihen erforderlich, die zunächst hier abgeleitet werden sollen.

1. \gesperrt{Ist $a_1, a_2, a_3, \ldots, a_n\ldots$ ein unbegrenztes System reeller (positiver, negativer oder auch verschwindender) Grössen, von dem wir voraussetzen, dass die Summe}

\weq{1}{\sigma_n = a_1+a_2+\cdots+a_n}

\gesperrt{für jedes beliebige $n$ dem absoluten Werthe nach unter einer endlichen Grenze $C$ bleibt, so ist die Reihe}

\weq{2}{S = \frac{a_1}{1^s}+\frac{a_2}{2^s}+\frac{a_3}{3^s}+\cdots}

\gesperrt{für jedes positive $s$ nicht nur convergent, sondern auch eine stetige Function von $s$.}

Um diesen Satz zu beweisen, zerlegen wir $S$ in der Weise:
\[
S = S_m+R_m,
\]
worin
\[
S_m = \frac{a_1}{1^s}+\frac{a_2}{2^s}+\cdots+\frac{a_{m-1}}{(m-1)^s},
\]
\[
R_m = \frac{a_m}{m^s}+\frac{a_{m+1}}{(m+1)^s}+\frac{a_{m+2}}{(m+2)^s}+\cdots
\]
Nun ist nach (1):
\[
a_m = \sigma_m-\sigma_{m-1}, \quad a_{m+1}=\sigma_{m+1}-\sigma_m, \quad a_{m+2}=\sigma_{m+2}-\sigma_{m+1}, \ldots
\]
und daher
\[
\begin{aligned}
R_m+\frac{\sigma_{m-1}}{m^s}
&= \sigma_m\left(\frac{1}{m^s}-\frac{1}{(m+1)^s}\right) \\
&\quad +\sigma_{m+1}\left(\frac{1}{(m+1)^s}-\frac{1}{(m+2)^s}\right)+\cdots
\end{aligned}
\]
Da nun nach der Voraussetzung $\sigma_{m-1}, \sigma_m, \sigma_{m+1}, \ldots$ zwischen endlichen Grenzen $\pm C$ eingeschlossen sind, so ist hiernach $R_m+\dfrac{\sigma_{m-1}}{m^s}$ zwischen den beiden Grenzen
\[
\pm C\left(\frac{1}{m^s}-\frac{1}{(m+1)^s}+\frac{1}{(m+1)^s}-\frac{1}{(m+2)^s}+\cdots\right) = \pm\frac{C}{m^s}
\]
%CONT>>
